% !TEX root = ../main.tex

\documentclass[../main.tex]{subfiles}

\graphicspath{{\subfix{../images/}}}

\begin{document}
	
	Актуальность исследования заключается в том, что инерциальная навигация является одним из ключевых
	методов определения положения и ориентации объектов в пространстве 
	без использования внешних источников сигналов. Это особенно важно для 
	автономных летательных аппаратов, функционирующих в условиях отсутствия или 
	недоступности спутниковых навигационных систем. Современные полётные 
	контроллеры, такие как Pixhawk, оснащены инерциальными измерительными 
	модулями (IMU~--- \textit{Inertial Measurement Unit}), позволяющими получать
	данные об ускорениях и угловых скоростях в реальном времени. Однако 
	обработка этих данных сопряжена с проблемами накопления ошибок интегрирования,
	что делает задачу повышения точности и устойчивости вычислений актуальной. 
	
    Новизна работы заключается в разработке программной системы, объединяющей 
	сбор сырых IMU-данных по протоколу MAVLink, их обработку с применением
	кватернионного интегрирования и методов компенсации дрейфа, а также 
	визуализацию результатов в реальном времени в интерактивной 3D-среде.

	Цель исследования: Разработка программной системы для считывания, обработки 
	и визуализации данных инерциальной навигации с полётного контроллера Pixhawk 
	в реальном времени.
	
	Объект исследования: Процесс инерциальной навигации автономных летательных 
	аппаратов на основе данных IMU.
	
	Предмет исследования: Методы обработки и интегрирования данных акселерометра
	и гироскопа, а также алгоритмы компенсации накопленных ошибок.
	
	Задачи работы:
	\begin{enumerate}[leftmargin=*]
		\item Реализовать модуль первичного считывания данных по протоколу
			MAVLink через последовательный порт с записью в CSV.
		\item Разработать многопоточную архитектуру приложения.
		\item Реализовать алгоритм калибровки и кватернионого интегрирования
			данных с гироскопа.
		\item Реализовать алгоритм двойного интегрирования акселерометра для
			оценки позиции.
		\item Создать OpenGL-визуализацию с 3D-моделью полетного контроллера,
		траекторией движения и HUD-панелью графиков датчиков.
		\item Провести анализ полученных результатов.
	\end{enumerate}

    Гипотеза исследования: Использование кватернионного
	интегрирования совместно с методами коррекции, такими как ZUPT, 
	позволит снизить накопление ошибок и повысить точность оценки 
	положения и ориентации объекта.

	Практическая значимость заключается в том, что разработанная система
	может быть использована для тестирования и отладки алгоритмов инерциальной
	навигации, а также в образовательных и исследовательских целях 
	при изучении работы IMU и систем управления беспилотными аппаратами.

\end{document}